\documentclass[12pt]{article}
\usepackage[utf8]{inputenc}
\usepackage[T2A]{fontenc}
\usepackage[russian]{babel}
\usepackage{amsmath, amssymb}
\usepackage{geometry}
\geometry{a4paper, margin=2cm}

\begin{document}

\section*{Анализ системы переписывания слов}

Дана SRS:
\[
\begin{cases}
aaa \to aaab,\\
aabb \to abab,\\
aa \to aba,\\
aaab \to ab,\\
bbb \to baaaa.
\end{cases}
\]

\subsection*{Завершимость.}
Заметим, что первое правило можно бесконечно применять к слову aaa: 
\[
aaa \to aaab \to aaabb \to aaabbb \to \ldots
\]
Отсюда получаем, что исходная система не является завершимой. Далее переориентируем правила в выбранном фундированном порядке: сначала сравниваем длину, а затем уже лексикографически.

\subsection*{Конфлюэнтность.}
Рассмотрим перекрытия для \(abab \to aabb,\; ababab\), применив это правило к разным подстрокам слова получаем \((aabbab,\; abaabb)\), нормализовав получаем \((aabbab,\; aaa)\), элементы пары не равны между собой, потому система не локально конфлюэнтна.

\subsection*{Пополняемость.}
Пополняя при помощи кода, получаем систему:
\begin{align*}
& aaab \to aaa, \\
& abab \to aabb, \\
& aba \to aa, \\
& aaab \to ab, \\
& baaaa \to bbb, \\
& aaa \to ab, \\
& bbbaab \to bbbaa, \\
& bbbab \to bbba, \\
& bbbb \to bbb, \\
& aaaaaaa \to aaa, \\
& aabbab \to aaa, \\
& aabb \to aab, \\
& aabba \to aaa, \\
& aabb \to aaaa, \\
& bbbbab \to bbba, \\
& baaa \to abaa, \\
& abb \to ab, \\
& aab \to aa, \\
& abbb \to ab, \\
& abba \to aa, \\
& bbbaa \to bbb, \\
& bbba \to bab, \\
& baab \to bbb, \\
& abbbaa \to ab, \\
& abbba \to aa, \\
& abbab \to aa, \\
& bbb \to baa, \\
& baab \to bbb, \\
& bbbaa \to bbb, \\
& bbba \to bab, \\
& aab \to aa, \\
& abbb \to ab, \\
& abbba \to aa, \\
& aab \to aa, \\
& abbb \to ab, \\
& bbbaa \to bbb, \\
& aab \to aa, \\
& aab \to aa, \\
& baa \to aa, \\
& baa \to aa, \\
& baa \to aa, \\
& baa \to aa, \\
& baa \to aa, \\
& baa \to aa, \\
& bab \to ab, \\
& baa \to aa, \\
& baa \to aa, \\
& baa \to aa, \\
& baa \to aa, \\
& baa \to aa, \\
& bbaa \to baa, \\
& bbab \to bab, \\
& baaa \to abaa.
\end{align*}
В ней все критически пары уже будут сводится к одной нормальной форме, потому уже пополненная система конфлюэнтна, а значит исходная пополнима.

\subsection*{Конечность.}
Получаем следующие нормальные формы: bba, aa, ab, bb, ba, a, b, $\varepsilon$. К остальным словам же будут применимы правила переписывания. Из-за конечного числа нормальных форм система конечна.

\subsection*{Минимальная система.}
Теперь будем делать левые части системы как можно меньше за счет остальных, пока не остановимся на моменте, когда это уже невозможно будет сделать, получаем:
\begin{align*}
& aba \to aa, \\
& aaa \to ab, \\
& abb \to ab, \\
& bbb \to baa, \\
& aab \to aa, \\
& bab \to ab, \\
& baa \to aa, \\
& baaa \to abaa.
\end{align*}

\subsection*{Мета-тестирование.}
Получаем 2 инварианта в 2 системах:
\begin{enumerate}
    \item Если в системе была буква a, то она никуда не пропадет,
    \item Четность числа букв a.
\end{enumerate}

\end{document}